\documentclass{article}
\usepackage{iclr2025_conference,times}
\input{math_commands.tex}
\usepackage{hyperref}
\usepackage{url}
\usepackage{booktabs}
\usepackage{xcolor}
\usepackage{amsmath}
\usepackage{amssymb}
\usepackage{graphicx}
\graphicspath{{../figures/}{figures/}}

\title{Online Shift Detection and Conformal Adaptation for Deployed Safety Classifiers}

\author{
Jun Wen Leong \\
\texttt{leongjunwen@gmail.com}
}

\newcommand{\ie}{\textit{i.e.}}
\newcommand{\eg}{\textit{e.g.}}

\iclrfinalcopy

\begin{document}

\maketitle
\lhead{Preprint. Under review.}

\begin{abstract}
Safety classifiers deployed in production operate under a stationarity assumption that fails silently: when input distributions drift, accuracy degrades with no error signal until ground-truth labels arrive. We present an online monitor that detects distributional shift in classifier scores via a sliding-window KS statistic with empirically calibrated alarm thresholds. In a pre-registered factorial evaluation (4 classifiers $\times$ 5 shift conditions $\times$ 20 seeds $\times$ 2 window sizes; 800 cells), the monitor achieves 86.6\% valid detection (mean latency 39.5 steps) across synthetic-onset, real-jailbreak, and adversarial regimes; a classifier $\times$ shift interaction ($\eta^2 = 0.185$) shows that monitoring must be tuned per classifier. Attempting to recover post-detection coverage via weighted conformal prediction exposes a failure mode: density-ratio estimation collapses for generative classifiers because logistic regression separates source from target perfectly in 3584--4096-dimensional embedding space, clipping all importance weights to zero; projecting to $\leq 32$ dimensions restores coverage. We then extend the framework to gradient-based evasion and give the first threat-model characterisation of score-disagreement monitoring as a canary. We falsify three assumptions: that architectural diversity drives the signal (false, $\eta^2 = 0.011$), that it is generic out-of-distribution detection (false, GCG-specific, $p < 10^{-12}$), and that an adaptive attacker can suppress it (false while the canary is confident). We derive the exact security boundary, a confidence-gated equilibrium at which a monitor-aware attacker stalls at gap $= 1/(2\lambda)$, and provide a calibration-free scan martingale achieving false-alarm rate $\leq 1\%$ across all classifiers with no per-model tuning.
\end{abstract}

%==============================================================================
\section{Introduction}
\label{sec:intro}
%==============================================================================

Safety classifiers are a critical defense layer between a language model and its users. When deployed at scale, they operate under a stationarity assumption: the distribution of inputs tomorrow will resemble the distribution on which the classifier was calibrated. This assumption can fail through adversarial adaptation~\citep{zou2023universal}, organic linguistic drift, multilingual code-switching, and compositional attacks that chain benign components into harmful sequences.

The failure mode is silent~\citep{rabanser2019failing}. A classifier whose accuracy has degraded from 95\% to 80\% produces no error signal unless ground-truth labels arrive. In production safety systems, labels typically do not arrive in real time. By the time periodic offline evaluation detects the problem, the classifier may have been making unreliable decisions for days or weeks. This silent-failure mode has recently been documented for embedding-based safety classifiers, where small representation drift collapses accuracy while confidence remains high~\citep{sahoo2026collapse}.

We propose an online monitor with empirically calibrated alarm thresholds that tells deployers when their safety classifier has moved out of distribution, before the shift accumulates further. The system watches the distribution of classifier scores via a sliding-window KS statistic with alarm thresholds calibrated via null simulation. When shift is detected, a conformal abstention layer reweights its prediction sets to preserve coverage under the new distribution.

We address three research questions:

\paragraph{RQ1 (Detection).} Can sequential statistical tests on classifier outputs detect distributional shift online, with controlled false alarm rates, across diverse shift types and classifier architectures?

\paragraph{RQ2 (Adaptation).} Does weighted conformal prediction recover coverage guarantees after shift detection, compared to unweighted conformal sets?

\paragraph{RQ3 (Factor Importance).} In a factorial evaluation design, what proportion of variance in detection latency is attributable to classifier choice, shift type, and their interaction?

Beyond drift detection, deployed classifiers face a second threat: gradient-based evasion~\citep{zou2023universal}. Fine-tuned safety classifiers are robust to template-based jailbreaks (1\% success rate in our evaluation; see \S\ref{sec:autodan}), making gradient attacks the remaining operational threat class. We extend the monitoring framework with two contributions:

\paragraph{RQ4 (Canary Detection).} Under what conditions does score disagreement between a targeted and un-targeted classifier detect gradient-based evasion, and when does it fail?

\paragraph{RQ5 (Adversarial Robustness).} Can an adaptive, monitor-aware attacker suppress the disagreement signal while evading the target classifier?

\paragraph{Contributions.} We contribute: (1)~a pre-registered factorial evaluation (800 cells) establishing detection baselines across 4 classifiers $\times$ 5 shift types; (2)~a calibration-free scan martingale replacing empirical threshold tuning (FAR spread reduced from $5\times$ to $<2\times$); (3)~the first adversarial characterisation of score-disagreement monitoring across 4 threat tiers, including a formal equilibrium bound (gap $= 1/(2\lambda)$, validated within 95\% CI); (4)~empirical falsification of three assumptions: architecture diversity as detection mechanism ($\eta^2 = 0.011$), reversed scaling in frontier models (a token-budget artifact), and monitorability as intrinsic classifier property ($r = 0.21$, $p = 0.70$ within-family); and (5)~a deployment blueprint across 35 frontier models with Wilson CI performance guarantees ($\geq 71\%$ detection at $< 1.5\%$ FPR; $N{=}1000$ FPR validation on 10 Pareto-dominant models).

%==============================================================================
\section{Related Work}
\label{sec:related}
%==============================================================================

\paragraph{Sequential testing and confidence sequences.} Classical sequential analysis~\citep{wald1945sequential} provides stopping rules with controlled error rates, but fixed-sample-size guarantees do not extend to continuous monitoring. \citet{howard2021timeuniform} establish the foundational theory of time-uniform confidence sequences; \citet{waudbysmith2024estimating} extend this to practical betting-based constructions: by constructing a wealth process that is a non-negative supermartingale under the null, Ville's inequality yields time-uniform coverage. The guarantee holds simultaneously at all time steps, not just at a pre-specified stopping time. We employ an adaptive betting strategy within this framework for the growing-window variant of our detector. For the sliding-window variant used in our factorial evaluation, the time-uniform guarantee does not hold; we instead calibrate alarm thresholds empirically via null simulation (\S\ref{sec:far-calibration}), which provides finite-horizon control at the cost of the anytime guarantee.

\paragraph{Two-sample testing on streams.} Kernel MMD~\citep{gretton2012kernel} is the standard nonparametric two-sample test for high-dimensional data, but its original formulation assumes fixed samples. We adapt it to the streaming setting by maintaining a sliding window of embeddings and comparing against frozen reference statistics. The bandwidth is fixed at calibration time following the standard median heuristic, preventing adaptation of kernel parameters to the detection window. This is closer to the block-based kernel two-sample tests of \citet{zaremba2013btests}, which reduce the cost of MMD estimation and provide a foundation for later online extensions, than to the original batch MMD, though we do not claim optimality of the kernel choice.

\paragraph{Conformal prediction under covariate shift.} Split conformal prediction~\citep{vovk2005algorithmic} provides distribution-free coverage guarantees under exchangeability. When the test distribution shifts, exchangeability breaks and coverage degrades. \citet{tibshirani2019conformal} restore coverage by reweighting the empirical distribution of calibration nonconformity scores with density ratios estimated from the covariate shift, provided those density ratios are known or accurately estimated. Our weighted-on-alarm mode implements their approach, triggered by the shift detector rather than assumed available from unlabeled target covariates. The density ratio is estimated via logistic regression on classifier embeddings, a lightweight approximation that avoids the instability of kernel density estimation in high dimensions. It is well established that density-ratio estimation degrades in high dimensions; dimensionality reduction before estimation is a known remedy~\citep{stojanov2019lowdim,sugiyama2011direct}. Our contribution is not the reduction itself but the diagnosis of a specific failure mode: in generative safety-classifier embeddings, logistic regression achieves perfect linear separability, driving all importance weights to the clip floor and eliminating data-driven reweighting (\S\ref{sec:conformal-results}). \citet{gibbs2021adaptive} extend conformal prediction to the fully online setting without exchangeability; we note this as future work. For conformal prediction applied specifically to classification, \citet{romano2020classification} provide adaptive prediction sets that maintain valid coverage while minimizing set size; our abstention criterion (flag when $|C(x)| \neq 1$) follows their framework.

\paragraph{Safety classifier monitoring.} Prior work on monitoring deployed classifiers focuses on performance estimation from unlabeled data~\citep{garg2022leveraging} or drift detection via population-level statistics~\citep{rabanser2019failing}, finding that performance depends strongly on shift type and representation. Several recent works couple shift detection with online adaptation: \citet{podkopaev2022tracking} track the risk of a deployed model and detect harmful distribution shifts using sequential testing (requiring labels upon prediction), and \citet{prinster2025watch} unify monitoring and adaptation within a single weighted-conformal-test-martingale framework that detects changepoints and adapts online to covariate shift. Conformal prediction has been applied to content moderation via conformalized estimates of annotation disagreement~\citep{villatecastillo2024moderation}. Our work differs from these general-purpose monitoring frameworks in scope and emphasis: we conduct a targeted empirical study of LLM safety classifiers under diverse shift conditions, surfacing failure modes that general monitoring does not, specifically the density-ratio collapse mechanism in high-dimensional generative embeddings, the architecture $\times$ shift-type crossover interaction, cross-classifier anomaly detection as a distributional canary for adversarial evasion, and the CS-vs-KS deployment trade-off at low contamination. Architecturally, we use a separate sliding-window detector that triggers a downstream weighted conformal layer (rather than a unified martingale), though we view the architectural choice as secondary to the empirical contributions. The factorial evaluation design (crossing classifiers with shift conditions) and its variance decomposition reveal interaction effects invisible to single-classifier studies; we are not aware of prior work that analyzes classifier $\times$ shift interactions this way for safety classifiers, though variance-decomposition analyses of prompt sensitivity have recently appeared for general LLM evaluation~\citep{romanou2026brittlebench}.


%==============================================================================
\section{Methods}
\label{sec:methods}
%==============================================================================

\subsection{System Architecture}
\label{sec:architecture}

The monitor observes a stream of classifier outputs $(x_t, s_t, \mathbf{r}_t)$: input text $x_t$, unsafe-class probability $s_t \in [0,1]$, and penultimate-layer representation $\mathbf{r}_t \in \mathbb{R}^d$. The primary detection channel in the factorial evaluation is the KS detector on classifier scores (\S\ref{sec:ks}). The system also includes an MMD detector on embeddings (\S\ref{sec:mmd}), evaluated post-factorial on a representative subset (\S\ref{sec:channel-comparison}).

\subsection{Reference Window Calibration}
\label{sec:reference}

Before monitoring begins, the system collects $n_{\text{ref}}$ records under the known in-distribution regime and freezes the reference CDF $\hat{F}_{\text{ref}}$ (sorted reference scores), against which the KS statistic is computed. The reference CDF is frozen at calibration time. No adaptive estimation occurs during monitoring; this is critical for the validity of the sequential tests.

\subsection{KS Detector}
\label{sec:ks}

The KS detector tracks whether the marginal distribution of classifier scores has changed. It maintains a sliding window of the most recent $w$ scores and computes the one-sample Kolmogorov--Smirnov statistic:
\begin{equation}
D_w = \sup_x |\hat{F}_w(x) - \hat{F}_{\text{ref}}(x)|
\end{equation}
If the input distribution shifts in a way that changes how the classifier scores inputs, the empirical CDF of recent scores diverges from the reference CDF. The KS statistic measures the maximum pointwise divergence, making it sensitive to any location, scale, or shape change in the score distribution. The sliding-window KS test is a standard concept-drift detector; the KSWIN method~\citep{raab2020reactive} compares two adaptive sub-windows of a stream. Our variant differs in that we compare a sliding window against a \emph{frozen} reference CDF rather than two adaptive sub-windows; the frozen reference eliminates adaptation drift but requires an initial calibration period (\S\ref{sec:reference}). We adopt this as a lightweight, well-understood detection channel rather than as a methodological contribution.

\subsection{MMD Detector}
\label{sec:mmd}

The system includes an MMD detector that tracks whether the geometry of classifier representations has changed. It maintains a sliding window of the most recent $w$ embedding vectors and computes the unbiased MMD$^2$ between the reference embeddings and the window:
\begin{equation}
\widehat{\text{MMD}}^2_u(X, Y) = \frac{1}{m(m-1)} \sum_{i \neq j} k(x_i, x_j) + \frac{1}{n(n-1)} \sum_{i \neq j} k(y_i, y_j) - \frac{2}{mn} \sum_{i,j} k(x_i, y_j)
\end{equation}
with Gaussian kernel $k(x, y) = \exp(-\|x - y\|^2 / 2\sigma^2)$. The MMD is zero if and only if the two distributions are identical (for characteristic kernels), making it sensitive to shifts that move the embedding mass even when the score marginal is preserved. The MMD threshold is calibrated via a 1000-permutation bootstrap null on the pooled reference embeddings, set at the $(1{-}\alpha)$ quantile. The factorial evaluation uses the KS detector only; the MMD channel is evaluated separately on a representative subset (\S\ref{sec:channel-comparison}).

\subsection{Confidence Sequences and Alarm Logic}
\label{sec:cs}

Each detector's statistic is wrapped in a confidence sequence (CS) to control false alarm rates over the monitoring horizon.

The guarantee we want: an alarm should fire only when the input distribution has genuinely shifted, with probability of false alarm bounded by $\alpha$ over the entire monitoring period.

What we use in practice: a sliding-window Hoeffding bound. For a window of $n$ bounded observations in $[a, b]$, the confidence interval around the empirical mean has half-width:
\begin{equation}
h = (b - a)\sqrt{\frac{\log(1/\alpha)}{2n}}
\end{equation}
This provides valid coverage for each individual window. However, it does not provide the time-uniform guarantee $\Pr(\exists t: \mu \notin [L_t, U_t]) \leq \alpha$. Over a long monitoring horizon, the probability of at least one false alarm exceeds $\alpha$.

We close the gap via empirical FAR calibration (\S\ref{sec:far-calibration}). Rather than relying on the theoretical bound alone, we simulate the null distribution by running $N_{\text{cal}} = 50$ negative control streams through the full pipeline and set the alarm threshold at the 97th percentile of the maximum observed statistic.

An alarm fires when the reference value exits the confidence interval. Alarms are suppressed during a warmup period of $w$ steps to allow the window to fill.

\subsection{Multiplicity Correction}
\label{sec:multiplicity}

The system architecture supports parallel KS and MMD detectors with \v{S}id\'{a}k correction ($\alpha_{\text{per}} = 1 - (1-\alpha)^{1/k}$) to control the family-wise error rate. Since the factorial evaluation uses the KS detector only ($k = 1$), no multiplicity correction is applied to the reported results.

\subsection{Statistical Methodology}
\label{sec:stats}

All statistical analyses use the following specifications: BCa bootstrap CIs on means (10,000 resamples, seed $= 42$, bias-corrected and accelerated method); CIs on rates via Wilson Score interval at 95\% confidence; CIs on coverage proportions via Clopper--Pearson exact interval; permutation tests (10,000 permutations for pairwise comparisons; 1,000 for ANOVA, sufficient given all $p < 0.001$); significance level $\alpha = 0.05$ throughout. Holm--Bonferroni correction is applied to the 8 highlighted pairwise comparisons discussed in \S\ref{sec:detection}: (1) decoder vs encoder on paraphrase, (2) encoder vs decoder on adversarial suffix, (3) DeBERTa adversarial vs paraphrase, (4) Llama Guard paraphrase vs adversarial, (5) window 100 vs 200 (paired), (6) Llama Guard $\times$ code-switch vs grand mean, (7) ShieldGemma paraphrase vs adversarial, (8) FAR DeBERTa vs Text-Moderation; all 8 survive correction at the family-wise $\alpha = 0.05$ level (weakest adjusted $p = 0.044$).

\subsection{Empirical FAR Calibration}
\label{sec:far-calibration}

We run $N_{\text{cal}} = 50$ negative control streams (reference data only, no shift injected) through the full detection pipeline. For each stream, we record the maximum KS statistic observed over the monitoring horizon. The alarm threshold is set at the $p$-th percentile of these maxima; in the factorial evaluation, $p = 97$. The empirical calibration accounts for the sliding-window correlation structure, the specific window size, and the stream length, factors that the theoretical Hoeffding bound treats conservatively.

\subsection{Scan Martingale}
\label{sec:martingale}

The empirical FAR calibration of \S\ref{sec:far-calibration} produces a $5\times$ spread across classifiers (Text-Moderation 2.0\% vs DeBERTa 9.5\%). We replace this with a conformal test martingale that provides anytime-valid FAR control without per-classifier threshold tuning.

For each incoming score $s_t$, we compute a two-sided conformal $p$-value against the frozen reference CDF: $p_t = 2 \min(F_{\text{ref}}^+(s_t), F_{\text{ref}}^-(s_t))$. Conformal $p$-values are clipped at $\delta = 1/(n_{\text{ref}}+1)$ to prevent log-wealth divergence under discrete score distributions where the test point may fall outside all reference values. The reference CDF is constructed from a held-out calibration set of 500 benign prompts, disjoint from the test prompts used for FAR evaluation; all thresholds are fixed before test-set evaluation. The betting function uses the power method:
\begin{equation}
\log W_t = \log W_{t-1} + \log \varepsilon + (\varepsilon - 1) \log p_t
\end{equation}
where $\varepsilon = 0.3$ is the betting parameter. We maintain $M = 50$ concurrent sub-martingales (one started per timestep), each tested at threshold $\log(M/\alpha)$ (union bound). An alarm fires when any sub-martingale exceeds this threshold. Results are validated at $M \in \{50, 100\}$; broader window sweeps are deferred (see \S\ref{sec:limitations}).

By Ville's inequality applied to each active sub-martingale and a union bound over the $M = 50$ active windows, $P(\exists t: \text{alarm} \mid H_0) \leq \alpha$ under exchangeability. This finite-horizon guarantee applies to the bounded-window construction ($M$ windows, tested horizon) and does not extend to arbitrarily long monitoring horizons with unlimited martingale restarts beyond this configuration. Empirically, the scan martingale achieves FAR $\leq 1\%$ across all four classifiers with no per-classifier null-threshold calibration (vs.\ 2--10\% under empirical KS calibration); $\varepsilon$ and $M$ remain deployment hyperparameters.

While the martingale guarantees uniform FAR without per-model calibration, detection power depends on $\varepsilon$ and $W$: these design choices trade Type~I guarantee for Type~II detection delay. $\varepsilon = 0.3$ achieves 100\% detection for encoders but 57--77\% for decoders with wide null distributions. A practitioner may increase $\varepsilon$ for decoder-heavy deployments at the cost of slower detection.

\paragraph{Honesty note on detection power.} Under instantaneous (step-onset) shift at full mixing, per-condition calibrated KS thresholds (97th-percentile tuned per shift type) match or exceed the scan martingale's detection rate. The scan martingale's contribution is \emph{operational}: exchangeability-valid FAR $\leq \alpha$ with no per-classifier null-threshold tuning ($\varepsilon$ and $M$ remain shared hyperparameters). Against gradual-onset adversaries who ramp shift slowly, the scan martingale has substantially higher power than the sliding-window KS (see AD1 below; 67\% vs 7\% at 20\% mixing with a 50-step ramp). If a practitioner invests in per-condition KS calibration \emph{and} shift onset is abrupt, KS is at least as sensitive; for gradual or ramped onset, the scan martingale is preferred.

\paragraph{Monitor-aware adversary (AD1).} The martingale's advantage emerges against adversaries who exploit the KS detector's window structure. We simulate a ramped-onset adversary that gradually introduces shifted samples over a 50-step ramp. Against DeBERTa with small per-sample shift (Figure~\ref{fig:ad1}):

\begin{table}[h]
\centering
\small
\caption{Detection rate under ramped-onset adversary (DeBERTa, 50-step ramp, small per-sample shift). KS uses a single global 97th-percentile threshold; scan martingale uses $\varepsilon=0.3$, $W=50$.}
\label{tab:ad1}
\begin{tabular}{lcccc}
\toprule
\textbf{Mixing rate} & \textbf{KS detection} & \textbf{Scan detection} & \textbf{Gap (pp)} \\
\midrule
15\% & 3\% & 37\% & +34 \\
20\% & 7\% & 67\% & +60 \\
25\% & 10\% & 87\% & +77 \\
30\% & 47\% & 100\% & +53 \\
\bottomrule
\end{tabular}
\vspace{2pt}
{\footnotesize\textit{Note:} KS threshold here is a single global null-distribution threshold (97th-percentile pooled across conditions). Per-condition calibrated KS substantially narrows the gap at full mixing (see Honesty Note above); the scan martingale's advantage is largest under ramped onset with a conservatively-set global threshold.}
\end{table}

\noindent At 20\% mixing with ramped onset, KS is effectively blind (7\% detection) while the scan martingale catches two-thirds of injections.

\begin{figure}[t]
\centering
\includegraphics[width=0.6\textwidth]{figures/fig_ad1_killer_chart.pdf}
\caption{Detection rate under ramped-onset adversary. The scan martingale dominates KS at low mixing rates where the per-window signal is too diffuse for a single KS threshold exceedance.}
\label{fig:ad1}
\end{figure}

\paragraph{Residual blind spot.} Against a constant sub-threshold adversary (sustained low-rate injection at $\leq$12\% mixing without ramp), neither KS nor the scan martingale reliably detects. The martingale's advantage is specific to ramped or accelerating onsets.

For long-horizon deployments where benign drift may violate exchangeability, we recommend a rolling reference window (gradual drift triggers alarms on encoders but not decoders; block-structured violations always trigger). In production environments with high temporal auto-correlation, exchangeability can be further preserved by sub-sampling the stream at session-level boundaries.

\subsection{Conformal Abstention Layer}
\label{sec:conformal}

Upon alarm, the system activates a conformal prediction layer that adapts decision thresholds to preserve a target error rate $\epsilon$ without requiring new labeled data.

\paragraph{Unweighted mode.} Standard split-conformal prediction~\citep{vovk2005algorithmic}. Given $n$ calibration examples with known labels: (1) compute nonconformity scores $\alpha_i = 1 - f(x_i)_{y_i}$; (2) set threshold $\hat{q}$ at the $\lceil(1-\epsilon)(n+1)\rceil/n$-th quantile; (3) at test time, include class $y$ in the prediction set if its nonconformity score $\leq \hat{q}$; (4) abstain when the prediction set contains more than one class or is empty. Under exchangeability, this guarantees $1-\epsilon$ coverage. Under covariate shift, exchangeability breaks and coverage degrades.

\paragraph{Weighted-on-alarm mode.} After the shift detector fires, we estimate density ratios $w_i = p_{\text{target}}(x_i) / p_{\text{source}}(x_i)$ via logistic regression on source vs.\ target embeddings~\citep{tibshirani2019conformal}. The ratios are clipped to $[1/C, C]$ with $C = 10$ for stability. The conformal quantile is recomputed as a weighted quantile:
\begin{equation}
\hat{q}_w = \inf\left\{q : \sum_{i: \alpha_i \leq q} \tilde{w}_i \geq 1 - \epsilon\right\}
\end{equation}
where $\tilde{w}_i = w_i / (\sum_j w_j + 1)$. This reweights the calibration scores to account for the covariate shift, restoring the coverage guarantee under the new distribution provided the density ratios are accurately estimated. The test-point contribution $w(X_{n+1})$ is set to 1.0 following \citet{tibshirani2019conformal}; when calibration weights collapse to the floor value of 0.1, this asymmetry shifts the effective quantile level by approximately 3 percentage points, producing the small residual recoveries ($\leq$+0.10) visible in Table~\ref{tab:conformal} at ESS$\approx$300.

\subsection{Variance Decomposition}
\label{sec:anova}

To quantify which experimental factors drive detection latency, we fit a two-way fixed-effects ANOVA:
\begin{equation}
\text{latency}_{ijk} = \mu + \alpha_i^C + \beta_j^S + (\alpha\beta)_{ij}^{CS} + \epsilon_{ijk}
\end{equation}
with factors classifier ($C$, 4 levels) and shift type ($S$, 5 levels). We report $\eta^2$ (proportion of total sum of squares) for each term. Bootstrap confidence intervals (1000 resamples, percentile method) quantify uncertainty in the main-effect estimates. The interaction term $(\alpha\beta)^{CS}$ captures classifier--shift pairings that are systematically easier or harder than predicted by the marginal effects alone.


%==============================================================================
\section{Experimental Setup}
\label{sec:setup}
%==============================================================================

\subsection{Classifiers}
\label{sec:classifiers}

\begin{table}[t]
\caption{Safety classifiers evaluated. Encoders fine-tuned on WildGuardMix (binary classification); decoders use original pre-trained weights.}
\label{tab:classifiers}
\centering
\small
\begin{tabular}{llccc}
\toprule
\textbf{Classifier} & \textbf{Architecture} & \textbf{Parameters} & \textbf{Embedding dim} & \textbf{Training data} \\
\midrule
DeBERTa-v3-large & Transformer encoder & 304M & 1024 & WildGuardMix \\
Text-Moderation (KoalaAI) & DeBERTa-v3-base & 86M & 768 & WildGuardMix$^\dagger$ \\
Llama Guard 3 & Decoder-only LLM & 8B & 4096 & Proprietary \\
ShieldGemma & Decoder-only LLM & 9B & 3584 & Proprietary \\
\bottomrule
\end{tabular}
\vspace{2pt}

{\footnotesize $^\dagger$Fine-tuned from \texttt{KoalaAI/Text-Moderation} (originally trained on proprietary data).}
\end{table}

The selection spans two architectural families (discriminative encoders and generative decoders) and two scales within each family, enabling analysis of both architecture and scale effects.

\subsection{Shift Conditions}
\label{sec:shifts}

\begin{table}[t]
\caption{Shift conditions. Each corpus is subsampled to 300 examples per factorial cell (temporal: 292 unique examples padded to 300 via repetition).}
\label{tab:shifts}
\centering
\small
\begin{tabular}{lll}
\toprule
\textbf{Condition} & \textbf{Mechanism} & \textbf{Threat model} \\
\midrule
Paraphrase & GPT-4o paraphrasing of harmful prompts & Organic rephrasing \\
Code-switch & Multilingual transliteration & Non-English users \\
Compositional & Multi-turn concatenation into long contexts & Context-window attacks \\
Temporal & Real jailbreaks from public red-team databases & Emerging harm categories \\
Adversarial suffix & GCG-optimized suffixes~\citep{zou2023universal} & Automated red-teaming \\
\bottomrule
\end{tabular}
\end{table}

The five conditions span the spectrum from naturalistic drift (paraphrase, temporal) to deliberate adversarial attack (GCG suffixes), with code-switch and compositional as intermediate cases.

\paragraph{Corpus validation.} We manually reviewed a sample from each corpus (50 examples for paraphrase and code-switch, 20 for compositional, temporal, and adversarial suffix). Paraphrase: approximately 14--20\% of examples became LLM refusals where the paraphrasing model declined to paraphrase harmful content, outputting safety responses instead. These contaminate the paraphrase condition with a refusal-text signal distinct from the intended semantic-preserving paraphrase shift; the monitor detects distributional change in the production stream regardless of mechanism, so detection latencies for paraphrase reflect a mixture of semantic shift and refusal artifacts and should be interpreted conservatively. A filtered ablation (removing the 9.4\% of examples identified as refusals, $n{=}5$ seeds $\times$ 2 classifiers) confirms this has negligible effect on detection: DeBERTa latency 38.0$\rightarrow$37.8 steps, Llama Guard 66.6$\rightarrow$60.8 steps, both 5/5 detected. Code-switch: all 50 examples were confirmed as authentic Singlish by a native speaker; 20--30\% became refusals, with the same caveat. Compositional: 20/20 correctly placed harmful content at the stated position (beginning/middle/end) within benign long-context wrappers; 100\% structural integrity. Temporal: 20/20 reviewed examples were genuine jailbreak prompts covering persona hijack, DAN-style injection, and NSFW extraction; zero false positives. The full temporal corpus (292 examples) draws from three public red-team databases: \texttt{lmsys/toxic-chat} (39\%), JailbreakBench (34\%), and ChatGPT-Jailbreak-Prompts (27\%). Adversarial suffix: 20/22 showed correct suffix concatenation with confirmed score flips (original $\geq 0.95 \rightarrow$ attacked $\leq 0.01$); one example was excluded post-validation (original score 0.002, already classified benign).

\subsection{Factorial Design}
\label{sec:factorial}

Full factorial: 4 classifiers $\times$ 5 shift conditions $\times$ 20 random seeds $\times$ 2 window sizes (100, 200) $= 800$ cells. Each cell runs a complete detection pipeline: reference window calibration $\rightarrow$ stream simulation with shift onset $\rightarrow$ alarm detection $\rightarrow$ latency measurement.

The factorial design, including all hyperparameters, was committed before execution (commit \texttt{be630f3}). Each cell includes a parallel negative control run (reference data only, no shift) to verify the alarm threshold does not fire on in-distribution data. A cell is marked as a valid detection only if: (1) detection latency $\geq 0$, and (2) the negative control does not alarm.

\paragraph{Compute.} Mac Studio (M2 Ultra, 192GB) for Llama Guard 3 and ShieldGemma; MacBook Pro (M3 Max, 128GB) for DeBERTa and Text-Moderation. Total wall-clock: ${\sim}120$ hours. All 800 cells completed without error.

\subsection{Deviations from Pre-Registration}
\label{sec:deviations}

The only scope reduction from the pre-registration is window sizes: 100 and 200 were executed; $w{=}500$ was dropped because it produces insufficient post-shift observations with 300 shifted examples. Seeds (20), ground-truth regimes (A, B, C), and all hyperparameters ($\alpha{=}0.05$, calibration percentile${=}97$, reference size${=}500$) match the pre-registration exactly.

\subsection{Reproducibility}
\label{sec:reproducibility}

Code, configurations, pre-registration document, and raw results are available at \url{https://github.com/junwenleong/safety-classifier-shift-monitor}. The pre-registration is anchored at commit \texttt{be630f3}; all hyperparameters are specified in version-controlled YAML files under \texttt{configs/}. Encoder fine-tuned checkpoints are not committed due to size; \texttt{scripts/finetune\_deberta.py} and \texttt{scripts/finetune\_text\_moderation.py} reproduce them from WildGuardMix with fixed seeds (${\sim}$2 hours each on Apple Silicon). A verification script (\texttt{scripts/verify\_paper\_numbers.py}) confirms the original factorial statistics against raw data (21 assertions, all passing); the extended conformal evaluation (Table~\ref{tab:conformal}) is verified against \texttt{results/conformal\_full.json}. Shifted corpora are generated offline with fixed seeds and committed before evaluation.


%==============================================================================
\section{Results}
\label{sec:results}
%==============================================================================

\subsection{RQ1: Detection Performance}
\label{sec:detection}

The system detects shift in 693 of 800 cells (86.6\% valid detection rate, 95\% Wilson CI [0.841, 0.888]), with empirical false alarm rates of 2--10\% across classifiers. ``Valid detection'' requires both a shift-stream alarm and a clean negative-control stream; raw TPR across alarming cells is 91.5\%, with the 5.8 pp gap attributable to false alarms on negative controls (Appendix~\ref{app:tpr-far}).

\begin{table}[t]
\caption{Mean detection latency (steps) with 95\% bootstrap CIs. $N{=}20$ seeds $\times$ 2 window sizes $= 40$ cells per combination; $n$ is valid detections.}
\label{tab:latency}
\centering
\small
\begin{tabular}{lccccc}
\toprule
\textbf{Classifier} & \textbf{Paraphrase} & \textbf{Code-switch} & \textbf{Compositional} & \textbf{Temporal} & \textbf{Adversarial} \\
\midrule
DeBERTa & 28.4 [26, 31] & 32.1 [30, 35] & 24.5 [22, 27] & 23.5 [22, 25] & 36.6 [33, 40] \\
Text-Mod. & 34.6 [32, 37] & 33.2 [31, 36] & 29.6 [28, 32] & 24.8 [23, 26] & 25.3 [23, 27] \\
Llama Guard & 69.4 [63, 76] & 93.4 [81, 107] & 47.0 [43, 51] & 42.1 [39, 45] & 27.8 [26, 30] \\
ShieldGemma & 85.0 [76, 94] & 81.8 [73, 90] & 43.9 [40, 48] & 27.1 [25, 29] & 26.8 [25, 29] \\
\bottomrule
\end{tabular}
\end{table}

The table reveals a crossover interaction. Paraphrase is easy for encoders (28--35 steps) but hard for decoders (69--85 steps). Adversarial suffix is the hardest condition for DeBERTa (36.6) but the easiest for Llama Guard (27.8). This crossover is not visible in any single-classifier study and motivates RQ3. The slowest cell is Llama Guard $\times$ code-switch (93.4 steps [81, 107]); the generation mechanism is largely invariant to surface-level script changes, requiring many shifted examples before the score distribution diverges detectably.

\begin{figure}[t]
\centering
\includegraphics[width=0.85\textwidth]{figures/fig1_latency_heatmap.pdf}
\caption{Detection latency heatmap (classifier $\times$ shift condition). Darker cells indicate slower detection. The crossover interaction is visible: encoders detect paraphrase fast but adversarial suffix slow; decoders show the opposite pattern.}
\label{fig:heatmap}
\end{figure}

\paragraph{Window size.} $w{=}100$ detects 7 steps faster on average (39.5 [37.0, 42.4] vs 45.4 [42.6, 48.3]; paired difference $-7.0$ [$-8.6$, $-5.3$], $n{=}314$ pairs) at the cost of slightly higher false alarm rates.

\paragraph{False alarm rates (95\% Wilson CIs).} Text-Moderation 2.0\% [0.8\%, 5.0\%] $<$ Llama Guard 3.0\% [1.4\%, 6.4\%] $<$ ShieldGemma 8.5\% [5.4\%, 13.2\%] $<$ DeBERTa 9.5\% [6.2\%, 14.4\%]. The $5\times$ spread reflects differences in score distribution shape under the null (Figure~\ref{fig:null-kde}): classifiers with wider score distributions (ShieldGemma $\sigma{=}0.14$, Llama Guard $\sigma{=}0.14$) produce more variable KS statistics, requiring higher calibrated thresholds and yielding higher residual FAR.

\begin{figure}[t]
\centering
\includegraphics[width=0.85\textwidth]{figures/fig_null_score_kde.pdf}
\caption{Null score distributions (in-distribution, $n{=}500$ per classifier). Text-Moderation and DeBERTa produce tightly concentrated scores near zero; Llama Guard and ShieldGemma show wider distributions with heavier tails, explaining the $5\times$ FAR spread across classifiers.}
\label{fig:null-kde}
\end{figure}

\paragraph{Failure analysis.} Of 107 invalid cells: 68 alarmed before shift onset (early alarm), 46 had dirty negative controls (false alarm), 7 exhibited both. Failures concentrate in DeBERTa $\times$ adversarial-suffix (13/40 invalid, 32.5\%) and Llama Guard $\times$ adversarial-suffix (11/40 invalid, 27.5\%). The DeBERTa failure is mechanistically explained by Regime C (\S\ref{sec:regime-c}): GCG suffixes push DeBERTa's scores into the safe region, producing no distributional signal. The Llama Guard failure reflects the opposite: adversarial tokens produce such a strong signal that the detector fires during calibration as well.

\subsection{RQ2: Conformal Adaptation}
\label{sec:conformal-results}

Evaluated across all four classifiers on three shift types (temporal, paraphrase, adversarial suffix), extending the initial temporal-only evaluation to test whether weighted conformal recovery is shift-type-dependent.

\begin{table}[t]
\caption{Conformal adaptation: coverage gap (unweighted) and recovery (weighted-on-alarm) across classifiers and shift types. ESS $= (\sum_i w_i)^2 / \sum_i w_i^2$; in this context ESS near the maximum ($n_{\text{cal}}{=}300$) indicates \emph{degenerate} uniformity (all calibration weights floor-clipped to $1/C$) rather than successful adaptation. Lower ESS reflects genuine density-ratio variation and more effective reweighting. $n{=}200$ post-shift examples per cell.}
\label{tab:conformal}
\centering
\small
\begin{tabular}{l ccc ccc ccc}
\toprule
& \multicolumn{3}{c}{\textbf{Temporal}} & \multicolumn{3}{c}{\textbf{Paraphrase}} & \multicolumn{3}{c}{\textbf{Adversarial suffix}} \\
\cmidrule(lr){2-4} \cmidrule(lr){5-7} \cmidrule(lr){8-10}
\textbf{Classifier} & Gap & Rec. & ESS & Gap & Rec. & ESS & Gap & Rec. & ESS \\
\midrule
DeBERTa & .085 & +.160 & 88 & .515 & +.390 & 46 & .325 & +.020 & 206 \\
Text-Mod. & .090 & +.020 & 300 & .515 & +.005 & 300 & $-.020$ & +.000 & 300 \\
Llama Guard & .350 & +.020 & 300 & .715 & +.015 & 300 & .260 & +.100 & 300 \\
ShieldGemma & .225 & +.075 & 300 & .725 & +.005 & 300 & .445 & +.060 & 300 \\
\bottomrule
\end{tabular}
\end{table}

The table reveals three findings. First, density-ratio collapse (ESS${\approx}$300, indicating uniform floor-clipped weights) is consistent across all three shift types for Text-Moderation, Llama Guard, and ShieldGemma. The collapse is a high-dimensional separability artifact, not a shift-specific phenomenon. Second, DeBERTa is the sole classifier where weighted conformal achieves meaningful recovery, but its effectiveness shows a gradient across shift types: paraphrase (ESS${=}$46, recovery${=}$+0.390) $>$ temporal (ESS${=}$88, recovery${=}$+0.160) $>$ adversarial suffix (ESS${=}$206, recovery${=}$+0.020). Third, DeBERTa on adversarial suffix shows near-total collapse (92.7\% of weights at floor, max weight 0.82) despite not reaching the binary collapse threshold; the GCG suffixes, optimized against DeBERTa specifically, push its embeddings into a region nearly perfectly separable from reference.

\paragraph{Mechanism.} The logistic regression density ratio estimator achieves perfect source/target separability in all collapsed cases, assigning $\hat{P}(\text{target} \mid x) \approx 0$ to every calibration point. All weights clip to the floor $1/C = 0.1$, eliminating data-driven reweighting. The residual $\leq$+0.10 recoveries observed at ESS$\approx$300 are explained by the test-point term: the implicit weight $w(X_{n+1})=1.0$ raises the effective quantile level from $90.3\%$ to $93.3\%$ at $n_{\text{cal}}=300$, $\varepsilon=0.1$: a formula artifact, not adaptation. ESS is the continuous predictor of recovery effectiveness: ESS${=}$46 yields +0.390 recovery; ESS${=}$88 yields +0.160; ESS${\geq}$200 yields ${\leq}$+0.020.

\paragraph{Dimensionality reduction as diagnostic.} To confirm the collapse is a dimensionality artifact, we applied PCA to 32 dimensions before density ratio estimation. This breaks the collapse for both generative classifiers on temporal shift: Llama Guard recovers 33.0 pp (coverage 0.555 $\rightarrow$ 0.885, ESS drops from 300 to 19.6), and ShieldGemma recovers 20.5 pp (0.690 $\rightarrow$ 0.895, ESS drops from 300 to 85.0). At 32 dimensions, 82--91\% of embedding variance is retained, but the logistic classifier can no longer achieve perfect separability. At 64 dimensions, ShieldGemma re-collapses (ESS${=}$300) while Llama Guard remains fixed (ESS${=}$135); at 128 dimensions, both re-collapse. The critical threshold is ${\leq}$32 dimensions for these embedding spaces. DeBERTa (1024-d, no baseline collapse) shows no degradation under PCA, confirming the reduction removes noise dimensions rather than safety-relevant signal. This result demonstrates that the collapse is driven by the high-dimensional separability identified by \citet{stojanov2019lowdim} in the general covariate-shift setting; the specific contribution here is diagnosing the failure mode in generative safety-classifier embeddings and showing it eliminates data-driven reweighting.

\subsection{RQ3: Variance Decomposition}
\label{sec:variance}

\begin{table}[t]
\caption{Two-way ANOVA on detection latency (693 valid detections). Main-effect 95\% CIs are percentile bootstrap (1000 resamples); interaction CI is from seed-clustered bootstrap (Appendix~\ref{app:anova-robustness}).}
\label{tab:anova}
\centering
\small
\begin{tabular}{lccc}
\toprule
\textbf{Factor} & $\boldsymbol{\eta^2}$ & \textbf{95\% CI} & \textbf{Permutation $p$} \\
\midrule
Classifier & 0.243 & [0.205, 0.291] & $< 0.001$ \\
Shift type & 0.237 & [0.193, 0.293] & $< 0.001$ \\
Classifier $\times$ Shift & 0.185 & [0.164, 0.223] & $< 0.001$ \\
Residual & 0.335 & -- & -- \\
\bottomrule
\end{tabular}
\end{table}

All three systematic factors are significant (1000 permutations, all $p < 0.001$). Classifier (24.3\%) and shift type (23.7\%) are the two largest systematic factors, with the interaction (18.5\%) close behind. The three factors contribute roughly equally; neither classifier choice nor shift type alone determines detection difficulty, and their interaction adds substantial additional variance.

\begin{figure}[t]
\centering
\includegraphics[width=0.7\textwidth]{figures/fig2_variance_decomposition.pdf}
\caption{Variance decomposition of detection latency. All three systematic factors contribute substantially (18--24\% each), with residual variance at 33.5\%.}
\label{fig:variance}
\end{figure}

The top interactions by magnitude: DeBERTa $\times$ adversarial-suffix ($+21.6$ steps above expected), ShieldGemma $\times$ paraphrase ($+19.9$), Llama Guard $\times$ code-switch ($+17.9$), Llama Guard $\times$ adversarial-suffix ($-16.0$). A monitoring system that sets thresholds based on classifier-level or shift-level averages will systematically under-alert on hard pairings and over-alert on easy ones.

\paragraph{Replication stability.} The initial $N{=}5$ estimate (200 cells) yielded $\eta^2 = 0.265$ for the interaction, which appeared to be the largest factor. At $N{=}20$ (800 cells), the interaction shrinks to 0.185 while the main effects grow (classifier: $0.196 \rightarrow 0.243$, shift: $0.217 \rightarrow 0.237$). The inflated $N{=}5$ interaction estimate is consistent with small-sample noise amplifying cell-level outliers. The qualitative finding (that the interaction is significant and substantial) replicates; the quantitative claim that it ``dominates'' does not.

\subsection{Robustness Across Ground-Truth Regimes}
\label{sec:regimes}

The preceding results use Regime A: synthetic shift onset at a known step. We evaluate two additional ground-truth regimes.

\paragraph{Regime B (Temporal split).} Real temporal jailbreaks from public red-team databases replace the synthetic shift corpus. The monitor receives no signal about when or whether shift occurs. Evaluated on DeBERTa and Llama Guard $\times$ temporal corpus $\times$ 5 seeds $\times$ 2 window sizes $= 20$ cells. The system detects shift in 17 of 20 cells (85\% detection rate, Wilson 95\% CI [64\%, 95\%]), with mean detection latency of 32.6 steps. Llama Guard achieves 10/10; DeBERTa achieves 7/10. Detection generalizes to naturally-occurring shift.

\paragraph{Regime C (Adversarial success).}
\label{sec:regime-c}
We filter GCG-optimized suffixes to the 22 examples that successfully flip DeBERTa's classification (one example excluded post-validation: original score 0.002, already classified benign), then stream them through all four classifiers. The result is a classifier-level demonstration of how the same adversarial inputs produce opposite distributional effects across architectures.

DeBERTa (target classifier): adversarial suffixes shift scores toward safe (post $\approx 0.01$), making attacked examples indistinguishable from the safe reference distribution. The monitor fails to detect in 38/40 seed-window combinations.

Llama Guard (non-target classifier): the same adversarial suffixes shift scores toward unsafe (post $= 0.78$), producing a massive upward shift ($+0.73$) that the monitor detects in 14/40 seed-window combinations.\footnote{Score direction diagnostics (post-shift mean scores and deltas) were computed at evaluation time via \texttt{scripts/check\_regime\_c\_direction.py} and are not reproduced by the verification script, which checks only detection counts. Full KS statistic time series for all four classifiers are stored in \texttt{results/regime\_c\_ks\_series.json} and can be reproduced without model inference.} A score direction analysis confirms all 14 detections are TOWARD\_UNSAFE: the GCG suffixes did not transfer their adversarial effect to Llama Guard. They made inputs appear more dangerous, not safer.

\begin{figure}[t]
\centering
\includegraphics[width=0.85\textwidth]{figures/fig_regime_c_ks_trajectories.pdf}
\caption{Regime C: KS statistic trajectories normalized by per-classifier threshold ($y{=}1.0$ is the alarm line). All classifiers hover near but below their thresholds on the 22-example adversarial corpus, explaining why detection is marginal. DeBERTa exceeds threshold briefly (false alarm before onset); the others remain sub-threshold.}
\label{fig:regime-c-ks}
\end{figure}

The correct interpretation is cross-classifier anomaly detection, not adversarial transfer: adversarial perturbations optimized against one classifier appear anomalous to architecturally different classifiers, enabling detection via score distribution shift even without attack transfer.

\paragraph{Corpus limitation.} The Regime C corpus comprises 22 successful adversarial examples; the identical Llama Guard score across examples (0.78) reflects a fixed adversarial token pattern rather than distributional variation. The 14 Llama Guard detections are replications of a single observation across seeds, not independent shift events. With window size 200 and calibrated thresholds, no classifier reliably detects shift from this corpus (Figure~\ref{fig:regime-c-ks}): all four KS trajectories hover near but below their respective thresholds, confirming the 22-example corpus is at the margin of detectability. The finding is robust and reproducible (full KS time series stored), but should be interpreted as a classifier-level result rather than a cell-level detection rate. A larger adversarial corpus with varied attack patterns would be needed to establish reliable cross-classifier detection.


%==============================================================================
\section{Discussion}
\label{sec:discussion}
%==============================================================================

\subsection{Per-Classifier Monitoring Profiles}
\label{sec:profiles}

The variance decomposition shows that classifier ($\eta^2 = 0.243$), shift type ($\eta^2 = 0.237$), and their interaction ($\eta^2 = 0.185$) all contribute substantially to detection latency. The interaction magnitude is large enough that a uniform monitoring configuration will systematically miscalibrate: under-alerting on hard pairings and over-alerting on easy ones.

Concretely: DeBERTa requires aggressive monitoring for adversarial suffixes ($+21.6$ steps above expected) but can afford relaxed thresholds for paraphrase. ShieldGemma requires the opposite profile: tight monitoring for paraphrase ($+19.9$ steps) but standard thresholds for adversarial inputs.

\subsection{Encoder vs.\ Decoder Safety Models}
\label{sec:encoder-decoder}

The factorial reveals two distinct architectural patterns that should not be conflated.

\paragraph{RQ1 (Detection latency): clean encoder/decoder split.} Discriminative models (DeBERTa, Text-Moderation) detect paraphrase and compositional shift fast (26--37 steps), while being slower on adversarial suffixes. Generative safety models (Llama Guard, ShieldGemma) show the opposite: slow on paraphrase (64--97 steps) and fast on adversarial suffixes (26--28 steps). This crossover interaction separates cleanly by architecture family in our four-model evaluation: surface-level rephrasing alters tokens that discriminative models attend to directly, while adversarial suffixes disrupt the generation distribution more visibly than a classification head.

\paragraph{RQ2 (Conformal recovery): not an encoder/decoder split.} The density-ratio collapse pattern does \emph{not} follow the same architectural divide. Text-Moderation (encoder, 86M, 768-d) collapses on all three shift types despite sharing training data (WildGuardMix) and architecture family (DeBERTa-v3) with the sole non-collapsing classifier. Only DeBERTa-v3-large (304M, 1024-d) avoids collapse.

We observe collapse despite same-distribution fine-tuning in the smaller encoder (Text-Moderation, 86M, 768-d, 9-class pretrained) but not in the larger encoder (DeBERTa-v3-large, 304M, 1024-d, binary). We hypothesize this reflects the multi-class pretraining objective of Text-Moderation (9 content categories: S, H, V, HR, SH, S3, H2, V2, OK), which shapes its embedding space to cleanly separate content categories, making it trivially easy for logistic regression to distinguish any two distributions in that space. However, we cannot distinguish this from scale effects (86M vs 304M, 768-d vs 1024-d) without a controlled ablation. The mechanism question remains open; what is empirically clear is that collapse is near-universal and training-distribution alignment alone does not prevent it.

\subsection{Detection Channel Comparison: KS vs.\ CS vs.\ MMD}
\label{sec:channel-comparison}

\paragraph{CS growing-window detector.} On a representative subset (4 classifiers $\times$ 3 shifts $\times$ 10 seeds $= 120$ cells), the growing-window confidence sequence achieves 120/120 detection (100\%) with 0/40 false alarms on reference-only streams (empirical FAR $= 0.0$). The observed result is consistent with the time-uniform guarantee of \citet{waudbysmith2024estimating}: no false alarm fired at any point during monitoring across all 40 null streams. However, detection latency is approximately $2\times$ that of KS (e.g., DeBERTa $\times$ paraphrase: CS 57.2 vs KS 25.1 steps).

The operational picture changes at low mixing proportions. At 30\% contamination (where only 30\% of post-onset traffic is shifted), the CS detector achieves 97\% detection (29/30, DeBERTa $\times$ paraphrase, $n{=}30$ seeds) while KS achieves only 43\% (13/30; Fisher exact $p < 0.0001$, Wilson CIs [0.83, 0.99] vs [0.27, 0.61], non-overlapping). The growing-window CS accumulates evidence across the full history rather than comparing a fixed sliding window, enabling detection of weaker shifts that fall below the per-window KS threshold. At 50\%+ mixing, both detect reliably but KS is faster.

This defines a deployment profile: \textbf{KS is preferred when shift signals are strong} (high mixing, abrupt onset, fast detection needed) and \textbf{CS is necessary when shift signals are weak} (low mixing, gradual contamination, reliable detection prioritized over speed). Real-world drift is rarely 100\% contamination; the CS advantage at low mixing is operationally significant.

\paragraph{Gradual drift and mixing sensitivity.} A ramp-rate sweep (DeBERTa $\times$ paraphrase, mixing ramps from 0\% to target over 50--200 steps, $n{=}10$ per condition) reveals that detection latency scales with ramp duration: 94 steps at 50-step ramp, 210 steps at 200-step ramp ($9/10$ detected). The mixing-level sweep at fixed 50-step ramp shows detection rate increasing monotonically: 43\% at 30\% mixing, 100\% at 50\%+, with latency decreasing from 94 to 64 steps as mixing increases from 50\% to 100\%.

\paragraph{MMD on embeddings.} The MMD detector on penultimate-layer embeddings (4 classifiers $\times$ 3 shifts $\times$ 10 seeds, calibrated via 1000-permutation bootstrap null at $\alpha{=}0.05$) detects 120/120 shifted streams with detection at the first possible step (latency $= w{=}100$ for all cells). Empirical FAR is controlled: DeBERTa 3.3\%, Text-Moderation 3.3\%, ShieldGemma 0.0\%. Llama Guard shows FAR $= 10\%$ (30 null windows; the wider embedding distribution produces less stable null estimates at this sample size).

MMD provides \emph{immediate binary detection} with no latency gradation across shift types or classifiers. This contrasts with KS, which shows meaningful latency variation (23--93 steps) reflecting the crossover interaction. The two detectors answer different operational questions: KS measures \emph{how quickly} the score distribution diverges, providing graded severity assessment; MMD measures \emph{whether} the embedding geometry has changed at all, providing a guaranteed backstop alarm. The two-channel architecture uses KS for graded alerts and MMD as a high-sensitivity binary alarm.

\subsection{Mechanistic Hypothesis}
\label{sec:mechanistic}

The crossover interaction is consistent with differences in null score distribution geometry. Discriminative classifiers (DeBERTa $\sigma{=}0.087$, Text-Moderation $\sigma{=}0.066$) produce tightly concentrated score distributions near zero; generative classifiers (Llama Guard $\sigma{=}0.144$, ShieldGemma $\sigma{=}0.141$) produce wider distributions with heavier tails. Across four classifiers, null score standard deviation correlates positively with mean detection latency ($r{=}0.97$, $p{=}0.032$, $n{=}4$). However, subsequent within-family evaluation using DeBERTa checkpoints at varying training epochs ($n = 6$ encoder variants, null-score $\sigma$ range 0.066--0.153) reveals $r = 0.21$, $p = 0.70$; the original correlation was an artifact of the encoder/decoder architectural gap, not an intrinsic monitorability property. Null-score geometry does not predict detection difficulty within an architecture family. See Appendix~\ref{app:monitorability} for the full falsification.

The pattern is shift-type-specific. For paraphrase, code-switch, compositional, and temporal shifts, wider null distributions predict slower detection ($r = 0.70$--$0.97$). For adversarial suffix, the correlation reverses ($r{=}{-}0.20$): wider-distribution classifiers detect \emph{faster}, producing the crossover.

To test whether embedding-space displacement drives the same pattern, we computed L2 centroid displacement between reference and shifted embeddings per classifier $\times$ shift condition. The overall correlation between displacement and detection latency is non-significant ($r{=}{-}0.09$, $p{=}0.78$), ruling out embedding geometry shift magnitude as the primary detection driver. Displacement partially tracks latency for paraphrase ($r{=}0.75$) but opposes it for adversarial suffix ($r{=}{-}0.38$), consistent with the hypothesis that the detection signal is mediated by score-boundary geometry rather than representation-space distance. These results are computed across $n{=}4$ classifiers and should be interpreted as a suggestive mechanistic pattern rather than a robust empirical confirmation.

\subsection{Limitations}
\label{sec:limitations}

\begin{itemize}
\item \textbf{PCA diagnostic validated on temporal shift only (primary experiment).} The PCA-32 dimensionality reduction that breaks density-ratio collapse has been demonstrated on temporal shift with 33 pp and 21 pp recovery. A secondary sweep on cached embeddings is consistent with ESS reduction generalizing to paraphrase shift (Llama Guard ESS$=$32 at dim$=$32; ShieldGemma ESS$=$28 at dim$=$32), breaking separability as expected. However, the secondary sweep uses a different calibration split (pooled across seeds rather than fresh inference), producing coverage at ceiling and precluding direct measurement of coverage recovery. The primary result remains the reported recovery magnitude.
\item \textbf{Residual variance.} 33.5\% of variance in detection latency is attributable to seed/noise. With 20 seeds per cell, the MDE is 13.9 steps at 80\% power.
\item \textbf{Binary classifiers only.} All four classifiers produce a single unsafe probability. Multi-category safety taxonomies may exhibit category-specific shift patterns invisible to a scalar score monitor.
\item \textbf{Canary analysis scope.} The adversarial robustness characterisation (\S\ref{sec:adversarial-robustness}) uses $n = 49$ GCG attacks on a single target (DeBERTa). Cross-target generalisation is untested.
\item \textbf{Joint-optimisation sample size.} Joint GCG results (\S\ref{sec:joint-gcg}) use $n = 20$ prompts at a single $\lambda = 2.0$ (11/20 joint flip, Wilson 95\% CI [34\%, 74\%]). Multi-$\lambda$ sweeps and larger-scale evaluations may reveal additional failure modes.
\item \textbf{Equilibrium derivation.} The stall condition (Appendix~\ref{app:derivation}) uses a continuous relaxation; GCG operates on discrete tokens. The match is empirical (within 95\% CI), not a formal convergence proof.
\item \textbf{Homogeneous negative controls.} Negative control streams draw exclusively from WildGuardMix unharmful examples; production streams with higher natural variance may require larger calibration sets.
\item \textbf{GCG suffix length.} A suffix-length sweep (10, 20, 40 tokens) on DeBERTa shows monotonically increasing attack success (70\%, 80\%, 90\% at $n = 10$; Fisher exact $p = 0.29$, underpowered). Transfer to frontier canaries is unaffected by suffix length: $\Delta = +0.26$ to $+0.30$ across all lengths (canary scores \emph{increase}, not decrease). Longer suffixes marginally improve local attack success without compromising non-transfer.
\item \textbf{Pre-registration deviation (window size 500).} The pre-registered factorial configuration included $w{=}500$; this was not executed because it produces insufficient post-shift observations with the 300-example shifted corpus. Window sensitivity analysis is restricted to $w \in \{100, 200\}$.
\item \textbf{Scan martingale window sweep.} The scan martingale is validated at $w \in \{50, 100\}$ (sub-martingale counts). A broader sensitivity analysis across $w \in \{25, 50, 75, 100, 150, 200\}$ was not performed; reported FAR and detection results hold only for the two validated configurations.
\item \textbf{Exchangeability in production traffic.} The scan martingale's FAR guarantee requires exchangeability under the null. Real production traffic is temporally correlated (multi-turn sessions, event-driven spikes). Deployers should sample across independent sessions rather than ingesting raw sequential streams to preserve the exchangeability assumption; benign domain drift that violates exchangeability will trigger alarms, which should be treated as triage signals for investigation rather than automatic enforcement.
\item \textbf{Fixed-$\lambda$ divergence-minimisation characterisation.} The equilibrium at gap $= 1/(2\lambda)$ characterises a fixed-$\lambda$ optimiser. A dynamic or annealing schedule for $\lambda$ could alter the optimisation trajectory; any stable convergence point still satisfies the instantaneous local gradient equilibrium, but we have not evaluated dynamic-$\lambda$ adversaries.
\item \textbf{Black-box attack scope.} The coordinate-ascent word-substitution evaluation (\S\ref{sec:llm-canary}) tests suffix-only black-box attacks (100 iterations). We did not evaluate stronger black-box methods such as genetic-algorithm prompt optimisation~\citep{liu2024autodan}, adaptive phrase-level rewriting, or multi-query prompt injection; robustness to these methods remains untested.
\end{itemize}

\subsection{Future Work}
\label{sec:future}

CUSUM or Bayesian online change-point detection for shift onset that accumulates over thousands of steps (below the 30\% mixing threshold identified in our ramp sweep). The equilibrium boundary (gap $= 1/(2\lambda)$) predicts a $\lambda$-dependent security frontier; a sweep across $\lambda$ values would confirm whether the stall point tracks the theory continuously. Cross-family joint optimisation (shared-tokeniser setting) remains untested due to the discrete-vocabulary incompatibility documented in \S\ref{sec:tokeniser}. Larger-scale adversarial evaluations ($n > 100$ prompts, multiple target classifiers) are needed to confirm the confidence-gating threshold and divergence-minimisation results. Truly online conformal prediction under arbitrary distribution shift~\citep{gibbs2021adaptive} would eliminate the batch requirement for density ratio estimation.

%==============================================================================
% v2 NEW SECTIONS: Canary Detection + Adversarial Robustness + LLM Canary
%==============================================================================
% =============================================================================
% NEW SECTIONS FOR arXiv v2 — Insert before Conclusion
% =============================================================================

%==============================================================================
\section{Canary Detection: Score Disagreement Under Targeted Evasion}
\label{sec:canary}
%==============================================================================

We expand the 22-example canary observation from \S\ref{sec:regime-c} into a full characterisation of score-disagreement monitoring. The target is DeBERTa-v3-large (the most effective encoder in our factorial evaluation); the canary classifiers are Text-Moderation, Llama Guard~3, and ShieldGemma. We evaluate on $n=100$ GCG-optimised prompts, of which 49 successfully flip DeBERTa below 0.5.

\paragraph{Note on terminology.} ``Steps'' in this section refer to GCG discrete token-search iterations on a single prompt (the attacker's optimisation budget), not production-stream timesteps monitored by the martingale (\S\ref{sec:cs}).

\subsection{Setup and Scoring Protocol}
\label{sec:canary-setup}

For each of the 49 successful GCG attacks against DeBERTa, we score the adversarial input on all four classifiers. We define \emph{divergence} for a classifier pair $(A, B)$ as $|f_B(x_{\text{adv}}) - f_A(x_{\text{adv}})|$, where $A$ is the target and $B$ is the canary. A canary \emph{detects} the attack when the divergence exceeds a null threshold calibrated at the 95th percentile of divergence under clean inputs.

Both GCG suffixes and random-token controls use exactly 20 tokens (\texttt{SUFFIX\_LEN=20}), controlling for perturbation magnitude. The divergence difference reported below is attributable to the optimisation structure of GCG, not to perturbation size.

\subsection{Attack-Specificity (CA6)}
\label{sec:ca6}

Score disagreement is not a generic out-of-distribution signal. We compare GCG-optimised suffixes against random-token suffixes of identical length (20 tokens) appended to the same prompts.

\paragraph{Result.} GCG attacks produce a detection rate of 75.5\% (37/49, Wilson 95\% CI [0.619, 0.854]) on the DeBERTa$\leftrightarrow$Llama Guard pair, versus 12.2\% (6/49, Wilson 95\% CI [0.057, 0.242]) for random-token controls. The mean divergence under GCG is 0.784 versus 0.213 under random tokens (Mann--Whitney $U$, $p = 2.37 \times 10^{-12}$).

The signal is attack-specific: gradient-optimised perturbations produce qualitatively different score-disagreement signatures than random noise of the same magnitude. This directly addresses the alert-fatigue problem --- the canary is silent under gibberish and triggers only under directed optimisation.

\subsection{Target-Specificity}
\label{sec:target-specificity}

We evaluate all 6 classifier pairs to test whether the divergence signal is specific to pairs involving the targeted classifier.

\paragraph{Result.} Three of six pairs show statistically significant GCG-vs-random differences (two-sided Mann--Whitney, Bonferroni-corrected $\alpha = 0.05/6 = 0.0083$): DeBERTa$\leftrightarrow$Text-Moderation ($p = 3.67 \times 10^{-17}$), DeBERTa$\leftrightarrow$Llama Guard ($p = 2.04 \times 10^{-12}$), and DeBERTa$\leftrightarrow$ShieldGemma ($p = 0.0023$; survives correction). The three non-significant pairs (Text-Moderation$\leftrightarrow$Llama Guard, $p = 0.320$; Text-Moderation$\leftrightarrow$ShieldGemma, $p = 0.715$; Llama Guard$\leftrightarrow$ShieldGemma, $p = 0.862$) all exclude the target classifier. The canary effect is target-specific: it manifests in pairs that include the attacked model.

\subsection{Architecture Is Not the Mechanism (CA4)}
\label{sec:ca4}

The prevailing assumption in multi-model defence literature is that architectural diversity fortifies ensembles against evasion. We test this by comparing within-family and cross-family divergence.

\paragraph{Result.} A one-way ANOVA on divergence with architecture type (within-family vs.\ cross-family) as the factor yields $\eta^2 = 0.011$, $F = 3.30$, $p = 0.070$. Architecture type explains 1.1\% of divergence variance. The highest mean divergence in our evaluation (0.899) occurs \emph{within}-family (DeBERTa$\leftrightarrow$Text-Moderation), not cross-family.

The mechanism is not architectural heterogeneity but weight-configuration divergence: any un-targeted classifier --- even a same-family checkpoint --- generates a robust divergence signature under single-target optimisation. Architecture diversity \emph{is} relevant for transfer robustness (\S\ref{sec:transfer}), but not for passive detection signal strength.

\subsection{$k$-Classifier Scaling}
\label{sec:k-scaling}

We evaluate detection rate as a function of the number of canary classifiers $k \in \{1, 2, 3\}$, using an OR rule (alarm if any canary detects).

\paragraph{Result.} The best single canary (Text-Moderation) achieves 93.9\% (46/49). The best $k{=}2$ combination (Text-Moderation + Llama Guard) achieves 98.0\% (48/49). Adding a third canary ($k{=}3$: all three) provides no additional detections (48/49 = 98.0\%). We recommend $k{=}2$ as the practical choice within our evaluated range $k \in \{1, 2, 3\}$: it achieves 98\% detection with no improvement from $k{=}3$.

\subsection{Template Attacks (AutoDAN)}
\label{sec:autodan}

Fine-tuned classifiers are robust to template-based jailbreaks. In our evaluation, AutoDAN-style template attacks achieve a 1\% flip rate against DeBERTa (the classifier has seen similar patterns during WildGuardMix training). This makes gradient-optimised attacks (GCG) the remaining operational threat class for fine-tuned safety classifiers.

\subsection{Tokeniser Fragmentation}
\label{sec:tokeniser}

Cross-family classifier pairs (encoder $\leftrightarrow$ decoder) use different tokenisers. We measure the tokenisation ratio: the number of tokens produced by the canary's tokeniser divided by the number produced by the target's tokeniser, for the same adversarial suffix.

\paragraph{Result.} The mean tokenisation ratio is $1.729 \pm 0.161$ ($n = 100$). A 20-token suffix optimised for DeBERTa's tokeniser maps to approximately 35 tokens under Llama Guard's tokeniser. This fragmentation explains why cross-family transfer is structurally suppressed: a single token substitution in the target's vocabulary shifts downstream subword boundaries in the canary's tokenisation, desynchronising the discrete search space (\S\ref{sec:cross-family-barrier}).

%==============================================================================
\section{Adversarial Robustness of Score-Disagreement Monitoring}
\label{sec:adversarial-robustness}
%==============================================================================

Having established that score disagreement detects single-target attacks (\S\ref{sec:canary}), we now ask: can an adaptive attacker suppress the disagreement signal? We characterise the security boundary across a 4-tier threat model.

\subsection{Threat Model}
\label{sec:threat-model}

\begin{table}[t]
\caption{Threat model tiers for score-disagreement monitoring. Detection rate is for the DeBERTa (target) + Text-Moderation/Llama Guard (canary) configuration.}
\label{tab:threat-tiers}
\centering
\small
\begin{tabular}{llcl}
\toprule
\textbf{Tier} & \textbf{Attacker knowledge} & \textbf{Detection rate} & \textbf{Deployment guidance} \\
\midrule
1 & Monitor-ignorant (single-target GCG) & 98\% (48/49) & $k{=}2$ canaries sufficient \\
2 & Knows canary exists (hopes for transfer) & 93\% (46/49) & Monitor canary confidence \\
3 & Has canary weights (joint optimisation) & 30\% (3/10) & Cross-family canary \\
4 & Stealth-aware (divergence minimisation) & 60\% (6/10) & Confidence gating \\
\bottomrule
\end{tabular}
\end{table}

Tier~1 is the baseline from \S\ref{sec:canary}. Tiers~2--4 represent progressively stronger adversaries. We evaluate Tiers~2--4 below.

\subsection{Transfer Analysis and Confidence Gating}
\label{sec:transfer}

A Tier~2 attacker optimises against the target alone but hopes the adversarial suffix accidentally fools the canary (passive transfer). We evaluate this on $n = 40$ prompts scored on same-family canaries (DeBERTa $\rightarrow$ Text-Moderation).

\paragraph{Confidence gating.} Transfer rate depends sharply on the canary's baseline confidence. When the canary's original score $f_B(x_{\text{orig}}) \geq 0.99$, the transfer rate is 4.3\% (1/23). When $f_B(x_{\text{orig}}) < 0.99$, the transfer rate is 76.5\% (13/17). The difference is significant (Mann--Whitney $p = 6.61 \times 10^{-7}$, $n = 40$).

\paragraph{Interpretation.} The detection threshold ($f_B \geq 0.5$) is an empirical operating point for the binary transfer/no-transfer decision. The gradient-paralysis mechanism (\S\ref{sec:divergence-min}) requires deeper saturation ($f_B \approx 1.0$, where $\|\nabla f_B\| \to 0$). These are distinct properties: detection relies on the canary holding its score above 0.5; gradient paralysis relies on the canary being in a flat, high-confidence basin.

\subsection{Joint Optimisation (Within-Family)}
\label{sec:joint-gcg}

A Tier~3 attacker has white-box access to both the target and the canary and optimises a combined loss $L = 0.5 \cdot L_A + 0.5 \cdot L_B$ using GCG with search width 256 and 50 steps.

\paragraph{Result.} Of 10 prompts evaluated, 7 achieve joint flip (both classifiers below 0.5), compared to 3/10 that transfer passively under single-target optimisation on these same prompts. Joint optimisation substantially improves the attacker's success rate within-family ($70\%$ vs.\ $30\%$ passive transfer).

\paragraph{Failure cases.} In 3/10 prompts, joint optimisation fails despite having white-box access to both models. Two of these (prompts 3 and 9: $f_A = 0.873$, $f_B = 0.998$ and $f_A = 0.916$, $f_B = 1.0$) exhibit gradient interference --- the combined gradient cannot simultaneously reduce both scores within the 50-step budget. One additional case (prompt 9) shows the canary holding at $f_B = 1.0$ while the target is stuck at $f_A = 0.916$, consistent with the confidence-gating mechanism.

\paragraph{Cross-family barrier.} The 1.73$\times$ tokeniser fragmentation ratio (\S\ref{sec:tokeniser}) structurally disrupts coordinate alignment for discrete token search. Because a single character substitution shifts downstream subword boundaries in Model~B (but not Model~A), the discrete search space becomes desynchronised --- a step that improves the joint loss for Model~A's tokenisation may worsen it for Model~B's. Cross-family joint optimisation remains structurally untested due to this discrete-vocabulary incompatibility.
\label{sec:cross-family-barrier}

\subsection{Divergence Minimisation: The Phase Transition}
\label{sec:divergence-min}

A Tier~4 attacker explicitly minimises score disagreement while attacking the target. We formalise this as:
\begin{equation}
\label{eq:div-min-loss}
L(\delta) = f_A(x + \delta) + \lambda \cdot (f_B(x + \delta) - f_A(x + \delta))^2
\end{equation}
where $\lambda = 2.0$ controls stealth pressure and Model~A is the target (DeBERTa epoch-1), Model~B is the canary (DeBERTa epoch-5).

\paragraph{Empirical result.} Of 10 prompts evaluated: 4 achieve stealth evasion (both models flip, divergence $< 0.5$) and 6 are blocked (Model~A remains above 0.5). In all 6 blocked cases, the canary is highly confident ($f_B \geq 0.987$). In 2 of 4 stealth cases, at least one model was originally uncertain ($f_B < 0.02$ or $f_A < 0.50$ at baseline); the remaining 2 stealth cases both have $f_B = 1.0$ at baseline, indicating that high canary confidence is necessary but not sufficient for blocking --- the target's loss landscape also contributes.

\paragraph{The stall condition.} The combined gradient of Eq.~\ref{eq:div-min-loss} is:
\begin{equation}
\label{eq:combined-gradient}
\nabla L = [1 - 2\lambda(f_B - f_A)] \cdot \nabla f_A + 2\lambda(f_B - f_A) \cdot \nabla f_B
\end{equation}

The coefficient $[1 - 2\lambda(f_B - f_A)]$ inverts sign when the gap $g = f_B - f_A$ exceeds $1/(2\lambda)$. At $\lambda = 2.0$, the critical gap is $1/(2 \times 2) = 0.250$. Beyond this point, the optimiser is driven to \emph{increase} $f_A$ --- the defence co-opts the attacker's gradient.

\paragraph{Empirical validation.} The 6 blocked cases stall at gap values $\{0.165, 0.242, 0.249, 0.250, 0.251, 0.251\}$, with mean 0.235 and bootstrap 95\% CI $[0.207, 0.251]$. The predicted critical value $0.250$ falls within this confidence interval. Five of six cases cluster at $0.242$--$0.251$; one outlier at $0.165$ (a prompt where Model~A's baseline was already lower) pulls the mean down.

\paragraph{Gradient-norm effect.} When the canary is deep in its confident basin ($f_B \approx 1.0$), $\|\nabla f_B\| \to 0$. The condition for a descent step to reduce $f_A$ requires (see Appendix~\ref{app:derivation}):
\begin{equation}
\cos(\theta) \geq \left(1 - \frac{1}{2\lambda(f_B - f_A)}\right) \cdot \frac{\|\nabla f_A\|}{\|\nabla f_B\|}
\end{equation}
As $\|\nabla f_B\| \to 0$, the right-hand side diverges. No angle $\theta$ satisfies the inequality --- the optimiser is completely paralysed. The defence is strongest precisely when the canary is most confident.

\paragraph{Active rejection.} In 3 of the 6 blocked cases, Model~A's score drifts \emph{upward} during optimisation (\eg, $0.476 \to 0.748 \to 0.758$ --- the model recovers toward the safe regime, reversing the attacker's progress). This active rejection is the coefficient inversion made visible: once $g > 1/(2\lambda)$, the combined loss penalises further reduction of $f_A$ more than it rewards it, creating defensive backpressure that pushes the target score back above the detection threshold.

\paragraph{Caveat.} This analysis uses a continuous relaxation. GCG operates on discrete tokens; the stall is empirical (within the 95\% CI of the predicted value), not a formal proof of impossibility. Higher $\lambda$ would block more aggressively but also constrain the attacker's loss landscape in ways that may prevent meaningful optimisation entirely.

\subsection{Deployment Recommendations}
\label{sec:deployment-recs}

The phase transition directly yields a measurable, input-time security indicator:

\begin{itemize}
\item \textbf{Confident canary} ($f_B \geq 0.99$): Trust the divergence signal. Transfer rate is 4.3\% and divergence-minimisation stalls at the predicted equilibrium. The input is protected.
\item \textbf{Uncertain canary} ($f_B < 0.5$): The divergence signal is unreliable. The attacker can succeed via joint-flip (70\% within-family). Route to human review regardless of observed divergence.
\item \textbf{$k = 2$ classifiers:} Deploy one same-family canary (highest detection sensitivity, 93.9\%) and one cross-family canary (structural transfer barrier via tokeniser fragmentation). This provides both sensitivity and robustness against adaptive attackers.
\end{itemize}

We do not claim this defence is absolute. A Tier~3 attacker with white-box access to a same-family canary succeeds 70\% of the time via joint optimisation (\S\ref{sec:joint-gcg}). The contribution is not an unbreakable defence but a quantitative characterisation of the exact security boundary --- where the defence holds, where it fails, and what measurable property ($f_B$) predicts which regime applies.


% =============================================================================
% APPENDIX C — Insert in Appendix section
% =============================================================================

\section*{Appendix C: Derivation of the Stall Condition}
\label{app:derivation}

\paragraph{Setup.} Consider the divergence-minimisation loss from Eq.~\ref{eq:div-min-loss}:
\[
L(\delta) = f_A(x + \delta) + \lambda \cdot (f_B(x + \delta) - f_A(x + \delta))^2
\]

Let $g = f_B - f_A$ denote the score gap. Expanding the gradient with respect to the perturbation $\delta$:
\begin{align}
\nabla_\delta L &= \nabla f_A + 2\lambda g \cdot (\nabla f_B - \nabla f_A) \nonumber \\
&= \nabla f_A + 2\lambda g \cdot \nabla f_B - 2\lambda g \cdot \nabla f_A \nonumber \\
&= (1 - 2\lambda g) \cdot \nabla f_A + 2\lambda g \cdot \nabla f_B
\label{eq:full-gradient}
\end{align}

\paragraph{Coefficient inversion.} The coefficient on $\nabla f_A$ is $(1 - 2\lambda g)$. This is positive when $g < 1/(2\lambda)$ and negative when $g > 1/(2\lambda)$. In the negative regime, a descent step on $L$ moves \emph{against} $\nabla f_A$ --- that is, it increases $f_A$ rather than decreasing it.

At $\lambda = 2.0$: the critical gap is $g^* = 1/(2 \times 2) = 0.250$.

\paragraph{Gradient-norm condition.} For a descent step $-\nabla L$ to reduce $f_A$, we need $\nabla L \cdot \nabla f_A \geq 0$ (so that moving against $\nabla L$ moves against $\nabla f_A$, reducing $f_A$). Equivalently, for the attack to succeed, we need:
\begin{align}
\nabla L \cdot \nabla f_A &= (1 - 2\lambda g)\|\nabla f_A\|^2 + 2\lambda g \cdot \nabla f_B \cdot \nabla f_A \geq 0
\end{align}

Let $\theta$ be the angle between $\nabla f_A$ and $\nabla f_B$. When $g > 1/(2\lambda)$:
\begin{align}
2\lambda g \cdot \|\nabla f_B\| \cdot \|\nabla f_A\| \cdot \cos\theta &\geq (2\lambda g - 1) \cdot \|\nabla f_A\|^2 \nonumber \\
\cos\theta &\geq \frac{2\lambda g - 1}{2\lambda g} \cdot \frac{\|\nabla f_A\|}{\|\nabla f_B\|}
\label{eq:cos-condition}
\end{align}

\paragraph{Paralysis in the confident basin.} When the canary is confident ($f_B \approx 1.0$), we empirically observe gradient saturation ($\|\nabla f_B\| < 10^{-3}$ across all blocked cases). Under this regime, the second term in Eq.~\ref{eq:gradient} vanishes and the combined gradient reduces to $\nabla \mathcal{L} \approx (1 - 2\lambda g)\nabla f_A$. Any gradient step changes $f_A$ by $\Delta f_A \propto -(1 - 2\lambda g)\|\nabla f_A\|^2$: when $g > 1/(2\lambda)$, this is positive (the step \emph{increases} $f_A$), and when $g < 1/(2\lambda)$, negative (the step decreases $f_A$). The critical gap $g^* = 1/(2\lambda)$ is therefore a stall point where the net coefficient vanishes, trapping the optimiser. This is consistent with our data: blocked cases converge toward $g^* = 0.250$ from both directions (higher baselines drift down, lower baselines drift up). Assuming $\|\nabla f_A\| > 0$ (the attacker still has a viable direction), the cosine condition (Eq.~\ref{eq:cos-condition}) additionally shows that the RHS diverges when $\|\nabla f_B\| \to 0$, reinforcing the blockade.

\paragraph{Empirical validation.} At $\lambda = 2.0$, the predicted stall gap is $g^* = 0.250$. In 6 blocked cases (where the canary held $f_B \geq 0.987$), the observed stall gaps are $\{0.165, 0.242, 0.249, 0.250, 0.251, 0.251\}$, mean $= 0.235$, bootstrap 95\% CI $[0.207, 0.251]$. The prediction $0.250$ lies within this interval. Excluding one outlier at $0.165$ (a prompt where discrete token substitution overshot the continuous boundary --- a known artifact of coordinate-wise search on non-smooth landscapes), the remaining five cluster at $0.242$--$0.251$ with mean $0.249$.

\paragraph{Limitations.} This derivation assumes continuous differentiability. GCG operates on discrete tokens via coordinate-wise search; the ``gradient'' is an approximation via first-order token-embedding differences projected onto the discrete vocabulary. The empirical match (within 95\% CI) suggests that this first-order approximation preserves the stall geometry, but the theoretical derivation serves as a predictive heuristic rather than a discrete impossibility proof. An attacker with access to qualitatively different optimisation methods (second-order, evolutionary, or RL-based search) might circumvent the first-order stall, though at substantially higher computational cost.ch) might circumvent the stall in ways that first-order analysis does not predict.


%==============================================================================
\section{Conclusion}
\label{sec:conclusion}
%==============================================================================

We presented an online monitoring system that detects distributional shift in deployed safety classifiers with 86.6\% detection rate across 800 pre-registered factorial cells, mean latency of 39.5 steps at window size 100, and false alarm rates of 2--10\%. The system generalizes across three ground-truth regimes: synthetic onset, real temporal jailbreaks, and adversarial success. Variance decomposition reveals that classifier, shift type, and their interaction all contribute substantially to detection difficulty, motivating per-classifier monitoring profiles. A growing-window confidence sequence provides 100\% detection with 0\% FAR and outperformed the sliding-window KS in our low-mixing experiments (97\% vs 43\% at 30\% contamination; $p < 0.0001$), while KS provides faster detection when signals are strong. MMD on embeddings provides a complementary binary alarm that fires immediately on all tested shifts, completing a three-channel architecture with different operational roles. Upon detection, weighted conformal prediction recovers up to 39 pp of lost coverage for DeBERTa across three shift types, but density-ratio estimation collapses for all other classifiers (ESS${\approx}$300, weights uniformly clipped to floor). PCA to 32 dimensions supports the diagnosis that this collapse is a dimensionality artifact, recovering 33 pp for Llama Guard and 21 pp for ShieldGemma on temporal shift, with ESS reduction consistent with generalizing to paraphrase shift. Code, configurations, and raw results are available at \url{https://github.com/junwenleong/safety-classifier-shift-monitor}.

Our adversarial characterisation establishes that score-disagreement monitoring is robust when the canary is confident: divergence-minimisation stalls at the predicted equilibrium (gap $= 1/(2\lambda) = 0.250$; of 14/20 blocked cases, 9 near-equilibrium cases cluster at mean $0.2499$, overall blocked-case mean $0.218$, median $0.250$), and white-box adversarial suffixes do not transfer to frontier canaries regardless of source architecture (encoder-to-frontier $\Delta = 0.000$ at $n{=}49$; decoder-to-frontier $|\Delta| \leq 0.010$ at $n{=}10$). The deployment rule is operationally simple: trust the canary signal only when the canary is confident; route uncertain inputs to human review.

%==============================================================================
\bibliography{references}
\bibliographystyle{iclr2025_conference}

%==============================================================================
\appendix
\section{Full Detection Rate Matrix}
\label{app:detection-rates}

\begin{table}[h]
\caption{Detection rate by classifier $\times$ shift condition ($N{=}40$ cells each).}
\centering
\small
\begin{tabular}{lccccc}
\toprule
\textbf{Classifier} & \textbf{Paraphrase} & \textbf{Code-switch} & \textbf{Compositional} & \textbf{Temporal} & \textbf{Adversarial} \\
\midrule
DeBERTa & 33/40 (82\%) & 35/40 (88\%) & 31/40 (78\%) & 35/40 (88\%) & 27/40 (68\%) \\
Text-Moderation & 34/40 (85\%) & 39/40 (98\%) & 37/40 (92\%) & 36/40 (90\%) & 34/40 (85\%) \\
Llama Guard & 37/40 (92\%) & 40/40 (100\%) & 38/40 (95\%) & 40/40 (100\%) & 29/40 (72\%) \\
ShieldGemma & 32/40 (80\%) & 35/40 (88\%) & 37/40 (92\%) & 32/40 (80\%) & 32/40 (80\%) \\
\bottomrule
\end{tabular}
\end{table}

\section{Replication: $N{=}5$ vs $N{=}20$ Comparison}
\label{app:replication}

\begin{table}[h]
\caption{Key statistics at $N{=}5$ (200 cells) vs $N{=}20$ (800 cells).}
\centering
\small
\begin{tabular}{lcc}
\toprule
\textbf{Statistic} & $\boldsymbol{N{=}5}$ & $\boldsymbol{N{=}20}$ \\
\midrule
Detection rate & 86.5\% [0.811, 0.906] & 86.6\% [0.841, 0.888] \\
Grand mean latency$^\dagger$ & 40.8 [37.5, 44.1] & 42.5 [40.6, 44.5] \\
$\eta^2$ Classifier & 0.196 & 0.243 \\
$\eta^2$ Shift type & 0.217 & 0.237 \\
$\eta^2$ Interaction & 0.265 & 0.185 \\
$\eta^2$ Residual & 0.322 & 0.335 \\
MDE (80\% power) & 22.4 steps & 13.9 steps \\
\bottomrule
\end{tabular}
\end{table}

{\footnotesize $^\dagger$Unweighted mean of classifier $\times$ shift cell means, pooled across both window sizes ($w{=}100$ and $w{=}200$). The abstract reports 39.5 steps, which is the mean over valid detections at $w{=}100$ only.}

The primary finding (86.5\% detection rate) replicates exactly at $N{=}20$. The variance decomposition shifts: the interaction term shrinks from 0.265 to 0.185 while main effects grow, consistent with small-sample noise inflating the $N{=}5$ interaction estimate. Statistical power improves from MDE $= 22.4$ to $13.9$ steps.

\section{Robustness of Variance Decomposition}
\label{app:anova-robustness}

Each seed contributes observations at two window sizes, introducing within-seed correlation. To verify that this does not inflate $\eta^2$ estimates, we conduct two sensitivity analyses.

\paragraph{Mixed-effects model.} A linear mixed model with seed as random intercept (REML estimation) yields random-effect variance $= 6.08$ vs.\ residual variance $= 242.72$. Seed explains only 2.4\% of residual variance, confirming the within-seed correlation is negligible.

\paragraph{Seed-clustered bootstrap.} We resample seeds (not individual observations) with replacement and recompute $\eta^2$ for each bootstrap sample ($B{=}2000$).

\begin{table}[h]
\caption{Seed-clustered bootstrap CIs for $\eta^2$ (2000 iterations, 20 seed clusters).}
\centering
\small
\begin{tabular}{lccc}
\toprule
\textbf{Factor} & $\boldsymbol{\eta^2}$ & \textbf{95\% CI (clustered)} & \textbf{Stable?} \\
\midrule
Classifier & 0.243 & [0.209, 0.285] & Yes \\
Shift type & 0.237 & [0.201, 0.286] & Yes \\
Interaction & 0.185 & [0.164, 0.223] & Yes \\
Residual & 0.335 & [0.253, 0.395] & Yes \\
\bottomrule
\end{tabular}
\end{table}

All CIs exclude zero with wide margins. The main effects move by $<$4\% and the interaction by $<$2\% relative to the point estimates, confirming that the qualitative conclusions are robust to the dependence structure.

\section{Detection Rate Decomposition}
\label{app:tpr-far}

The ``valid detection'' metric reported in \S\ref{sec:detection} requires both successful detection (latency $\geq 0$) \emph{and} a clean negative control (paired null stream does not alarm). Table~\ref{tab:tpr-far} separates these components.

\begin{table}[h]
\caption{Decomposition of detection metrics. Raw TPR = proportion with latency $\geq 0$; Control FAR = proportion with dirty negative control; Valid = both conditions met.}
\label{tab:tpr-far}
\centering
\small
\begin{tabular}{lccc}
\toprule
\textbf{Classifier} & \textbf{Raw TPR} & \textbf{Control FAR} & \textbf{Valid Detection} \\
\midrule
DeBERTa & 87.0\% & 9.5\% & 80.5\% \\
Text-Moderation & 91.5\% & 2.0\% & 90.0\% \\
Llama Guard & 95.0\% & 3.0\% & 92.0\% \\
ShieldGemma & 92.5\% & 8.5\% & 84.0\% \\
\midrule
\textbf{Overall} & \textbf{91.5\%} & \textbf{5.8\%} & \textbf{86.6\%} \\
\bottomrule
\end{tabular}
\end{table}

The 5-percentage-point gap between raw TPR (91.5\%) and valid detection (86.6\%) is fully explained by the 46/800 cells with dirty negative controls (Control FAR = 5.8\%). Detection sensitivity is high across all classifiers; the valid-detection metric is conservative because it penalizes detectors that are also over-sensitive on in-distribution data.

\section{Falsified: Monitorability Is Not an Intrinsic Classifier Property}
\label{app:monitorability}

The mechanistic hypothesis of \S\ref{sec:mechanistic} reported $r = 0.97$ ($p = 0.032$, $n = 4$) between null-score standard deviation and mean detection latency. We test whether this correlation holds within a single architecture family by fine-tuning DeBERTa-v3-large at training epochs \{1, 3, 5, 10\}, producing four encoder variants with null-score $\sigma \in [0.109, 0.153]$.

Combined with the two original encoders (DeBERTa $\sigma = 0.087$, Text-Moderation $\sigma = 0.066$), we obtain $n = 6$ within-family data points. The Pearson correlation is $r = 0.21$ ($p = 0.70$). Including the two decoders ($n = 8$ total): $r = 0.42$ ($p = 0.30$). Neither meets the pre-registered threshold of $r > 0.6$, $p < 0.05$ at $n \geq 6$.

\paragraph{Diagnosis.} The original $r = 0.97$ was driven by a two-cluster structure: decoders (high $\sigma \approx 0.14$, high latency $\approx 60$--85) vs.\ encoders (low $\sigma \approx 0.07$--0.09, low latency $\approx 12$--25). Within-family, no relationship exists (epoch~3 has $\sigma = 0.144$ and latency 13; epoch~5 has $\sigma = 0.153$ and latency 24). Monitorability is not an intrinsic scalar property of score-distribution geometry.

\section{Derivation of the Stall Condition}
\label{app:derivation}

\noindent\fbox{\parbox{0.96\textwidth}{%
\textbf{Proposition 1} (Divergence-Minimisation Equilibrium).
Let $f_A, f_B$ be differentiable classifiers with scores in $[0,1]$, and let an attacker minimise $L(\delta) = f_A(x+\delta) + \lambda(f_B(x+\delta) - f_A(x+\delta))^2$. If the canary is confident ($f_B \approx 1$, $\|\nabla f_B\| \to 0$), the score gap $g = f_B - f_A$ stalls at $g^* = 1/(2\lambda)$. At this point, the gradient coefficient on $\nabla f_A$ inverts sign, and further descent on $L$ \emph{increases} $f_A$.

\smallskip\noindent\textbf{Experimental validation:} At $\lambda = 2.0$, predicted $g^* = 0.250$. Observed across $n = 20$ prompts: 14 blocked, of which 9 cluster within $\pm 0.05$ of the prediction with mean gap $= 0.2499$. The remaining 5 stall earlier (the attack stalls before the equilibrium, making the defence even safer).%
}}\medskip

Consider the divergence-minimisation loss:
\[
L(\delta) = f_A(x + \delta) + \lambda \cdot (f_B(x + \delta) - f_A(x + \delta))^2
\]

Let $g = f_B - f_A$. The gradient with respect to the perturbation $\delta$:
\begin{align}
\nabla_\delta L &= (1 - 2\lambda g) \cdot \nabla f_A + 2\lambda g \cdot \nabla f_B
\end{align}

The coefficient on $\nabla f_A$ is $(1 - 2\lambda g)$. This inverts sign when $g > 1/(2\lambda)$. At $\lambda = 2.0$, the critical gap is $g^* = 0.250$. Beyond this point, a descent step on $L$ moves against $\nabla f_A$, increasing $f_A$ rather than decreasing it. When the canary is confident ($f_B \approx 1.0$, $\|\nabla f_B\| \to 0$), the combined gradient reduces to $\nabla L \approx (1 - 2\lambda g)\nabla f_A$, and the stall point $g^* = 1/(2\lambda)$ traps the optimiser. Empirically, across $n = 20$ prompts, 14 are blocked and stall at gap values with mean $0.218$ and median $0.250$; 9 of 14 cluster within $\pm 0.05$ of the critical value with mean gap $= 0.2499$. The remaining 5 stall earlier (the attack stalls before the equilibrium, making the defence even stronger). The prediction $0.250$ falls within the bootstrap 95\% CI on the 9 near-equilibrium cases.

\end{document}
